% -------------------------------------------------------------------------------------
\section{SCET 证明的前提与地位}\label{sec:scet-status}

\begin{remark}[SCET 证明的前提]\label{rem:scet-assumptions}
SCET 中的全阶因子化“证明”（例如\cref{thm:scet-thrust}）依赖以下前提，每一条都须在具体过程中单独论证：
(1) \emph{模式完备}：所选模式（硬、共线、软或超软，必要时 Glauber）再现 QCD 在给定观测量下的全部领头红外区域。这\emph{就是}第~4~章的区域分析（\cref{thm:ir-leading-regions}）；SCET 本身不能证明它，而在 Minkowski 运动学中区域方法本身也没有一般证明（\cref{prop:ir-regions-minkowski}，[P]）。
(2) \emph{零箱减除的自洽}（\cref{def:scet-zero-bin}）。
(3) \emph{Glauber 模式相消或不出现}：对只有出射共线粒子的 $e^+e^-$ 过程，Glauber 区不形成钳夹；对强子—强子碰撞，这是证明中最困难的一步（\cref{sec:scet-glauber}）。
(4) \emph{领头幂次}：退耦与算符基都只在领头幂次成立；在次领头幂次，卷积可以出现端点发散，因子化的形式本身需要修改（第~21~章）。
前三条与图形方法（第~9~章）的三处困难一一对应：
\begin{center}\small
\renewcommand{\arraystretch}{1.1}
\begin{tabular}{@{}lll@{}}
\toprule
SCET 的前提 & 图形方法的困难 & 本书中的编号假设 \\
\midrule
模式完备 & 区域完备与幂次计数 & \cref{asm:intro-bridge}(Br1)，\cref{asm:appN-envelope} \\
零箱减除的自洽 & 减除的组合学 & \cref{asm:appN-collins} \\
Glauber 拉氏量与快度正规子 & Glauber 区的围道变形 & \cref{asm:appO-mixing,asm:appO-deformation,asm:appO-rewritten-cs} \\
\bottomrule
\end{tabular}
\end{center}
在这些前提下，图形证明与 SCET 证明的内容相同：二者都以领头区域为出发点，前者用 Ward 恒等式与减除，后者用规范对称性、场重定义与零箱减除；两者的全阶陈述都是 [A]。SCET 证明并不比 Collins 的图形证明更严格，二者互补：SCET 把区域的组合学变成拉氏量的对称性，并把重求和变成 RG 方程的求解；图形方法把每个前提落实到具体的被积函数与围道上。Collins 写道：“I have not been able to penetrate the SCET literature to properly understand its rationale”（我未能深入 SCET 文献以充分理解其逻辑依据）\cite[第~15.9~节]{Collins2011}；上表正是为了在两种语言之间建立逐条的对应。两种语言也可以混合使用：Feige--Schwartz 把钳夹面分析、振幅方法与有效场论技巧结合，给出了振幅在领头幂次上的硬—软—共线因子化的全阶论证（据其摘要，与正规子和规范的选择无关）\cite{Feige:2013zla,FeigeSchwartz2014}；按 Rothstein--Stewart 的评述，对末态粒子的产生，该论证以与 CSS 类似的逻辑避开了振幅层次的 Glauber 区\cite[第~10.1~节]{RothsteinStewart2016}。它的输入同样是钳夹面的结构，在本书的分级中与其余 SCET 论证同属 [A]。第~24~章把前提 (1) 与 (3) 写成前沿章节使用的编号假设（\cref{asm:front-scet-modes}）。
\end{remark}

\begin{tcolorbox}[enhanced,breakable,title=常见误解,fonttitle=\bfseries,
  colback=gray!3!white,colframe=gray!55!black,boxrule=0.6pt,arc=2pt,
  left=5pt,right=5pt,top=3pt,bottom=3pt,before skip=6pt,after skip=6pt]
\textbf{“SCET 独立于区域分析证明了因子化。”}不对：SCET 预设了模式的内容——那正是区域分析——还需要零箱减除并须处理 Glauber 模式；它的全阶陈述与图形论证一样是 [A]（\cref{rem:scet-assumptions}）。\par
\textbf{“区域方法是定理，因此模式的选择不成问题。”}欧氏的大动量、大质量极限有证明\cite{Smirnov:1990rz}；Minkowski 极限没有一般证明（\cref{prop:ir-regions-minkowski}），可能需要解析正规子，且据该文摘要，“区域等于 Newton 多面体的下表面”在 Minkowski 运动学中不成立\cite{Gardi:2024axt}；对强子对撞中的喷注间隙截面，区域分析在三圈识别出了微扰的活动—活动 Glauber 交换贡献\cite{Becher:2024kmk}——这类不易被识别的“隐藏”区域见第~24~章的\cref{prop:front-hidden-regions}（[V]）。SCET 的模式须由钳夹分析论证。\par
\textbf{“零箱减除自动消除了重复计算。”}只有在每个共线因子中都自洽地实施时才如此；它与软减除的等价性是在用维数正规化处理红外发散时于低阶显式验证的（\cref{def:scet-zero-bin}）。
\end{tcolorbox}

\paragraph{本章小结。}
树图层次积掉 $\phi_{\bar n}$ 与离壳模式是精确的，$W_n$ 由程函因子按次序求和得到；领头阶拉氏量的规范对称性与 BPS 场重定义使超软胶子只出现在算符的 $Y_n$ 中（[R]）。SCET$_{\rm II}$ 的快度发散有三种处理，它们的等价性在低阶显式验证；Glauber 模式由 Rothstein--Stewart 的算符描述，而其相消没有完整证明（二者均为 [A]）。推力分布的因子化是 [A]，其单圈因子、RG 一致性与对全 QCD 奇异项的再现是 [R]，LL 指数作为 QCD 陈述继承 [A]。这些工具用于第~16、18、20、21、24~章。
